% answers.tex —— 2007 年·北京卷·文科 · 参考答案与详细解析片段
% 由 tools/generate_projects.py 依来源总账自 2007/2007_数学真题_北京文_试题.md + 2007/2007_数学真题_北京文_解析正文.md + 2007/2007_数学真题_北京文_答案要点.md 生成，勿手改（改源 MD 后重跑）


\begin{gktip}
固定候选 PDF 仅含题面，没有答案、解析或评分标准。以下为编者依据固定题面撰写的推导，不代表官方答案或官方解析；来源发布者、版次、复制链与官方身份未核，\texttt{official_answer=null}。Q18 的独立性条件未明，按答案要点保持 \texttt{U/PENDING}。
\end{gktip}

\gkpart{一、选择题}

\begin{answer}
\gksolhead{1}{}

\textbf{题目}：

已知 $\cos\theta\cdot\tan\theta<0$，那么角 $\theta$ 是（　）

\begin{gkopts}(2)
\task （A）第一或第二象限角
\task （B）第二或第三象限角
\task （C）第三或第四象限角
\task （D）第一或第四象限角
\end{gkopts}

\Needspace{4\baselineskip}
\textbf{解析}

原式有意义时 $\cos\theta\ne0$，并且 $\cos\theta\tan\theta=\sin\theta$。因此条件化为 $\sin\theta<0$，即角在第三或第四象限。

\gkanswer{答案：C}

\end{answer}
\begin{answer}
\gksolhead{2}{}

\textbf{题目}：

函数 $f(x)=3^x\ (0<x\leq 2)$ 的反函数的定义域为（　）

\begin{gkopts}(4)
\task （A）$(0,+\infty)$
\task （B）$(1,9]$
\task （C）$(0,1)$
\task （D）$[9,+\infty)$
\end{gkopts}

\Needspace{4\baselineskip}
\textbf{解析}

当 $0<x\leq2$ 时，$3^x$ 的值域为 $(1,9]$。反函数的定义域等于原函数的值域。

\gkanswer{答案：B，$(1,9]$}

\end{answer}
\begin{answer}
\gksolhead{3}{}

\textbf{题目}：

函数 $f(x)=\sin 2x-\cos 2x$ 的最小正周期是（　）

\begin{gkopts}(2)
\task （A）$\dfrac{\pi}{2}$
\task （B）$\pi$
\task （C）$2\pi$
\task （D）$4\pi$
\end{gkopts}

\Needspace{4\baselineskip}
\textbf{解析}

$\sin2x-\cos2x=\sqrt2\sin(2x-\pi/4)$，角频率为 $2$，故最小正周期为 $2\pi/2=\pi$。

\gkanswer{答案：B，$\pi$}

\end{answer}
\begin{answer}
\gksolhead{4}{}

\textbf{题目}：

椭圆 $\dfrac{x^2}{a^2}+\dfrac{y^2}{b^2}=1\ (a>b>0)$ 的焦点为 $F_1,F_2$，两条准线与 $x$ 轴的交点分别为 $M,N$。若 $|MN|\leq 2|F_1F_2|$，则该椭圆离心率的取值范围是（　）

\begin{gkopts}(2)
\task （A）$\left(0,\dfrac{1}{2}\right]$
\task （B）$\left(0,\dfrac{\sqrt{2}}{2}\right]$
\task （C）$\left[\dfrac{1}{2},1\right)$
\task （D）$\left[\dfrac{\sqrt{2}}{2},1\right)$
\end{gkopts}

\Needspace{4\baselineskip}
\textbf{解析}

椭圆离心率 $e=c/a$，两条准线与 $x$ 轴交点的距离为 $|MN|=2a/e$，两焦点间距为 $2c$。题设等价于 $2a/e\leq4c$，即 $1\leq2e^2$。再结合 $0<e<1$ 得 $\sqrt2/2\leq e<1$。

\gkanswer{答案：D，$\left[\dfrac{\sqrt2}{2},1\right)$}

\end{answer}
\begin{answer}
\gksolhead{5}{}

\textbf{题目}：

某城市的汽车牌照号码由 2 个英文字母后接 4 个数字组成，其中 4 个数字互不相同的牌照号码共有（　）

\begin{gkopts}(2)
\task （A）$(C_{26}^{1})^2A_{10}^{4}$ 个
\task （B）$A_{26}^{2}A_{10}^{4}$ 个
\task （C）$(C_{26}^{1})^2 10^4$ 个
\task （D）$A_{26}^{2}10^4$ 个
\end{gkopts}

\Needspace{4\baselineskip}
\textbf{解析}

两个字母位置分别有 $26$ 种选择，题面没有要求字母互不相同；四个数字互不相同，有 $10\cdot9\cdot8\cdot7$ 种。因此总数为 $26^2A_{10}^{4}$。

\gkanswer{答案：A，$(C_{26}^{1})^2A_{10}^{4}$}

\end{answer}
\begin{answer}
\gksolhead{6}{}

\textbf{题目}：

\Needspace{5\baselineskip}
若不等式组

\nopagebreak[4]
\[
\begin{cases}
x-y+5\geq 0,\\
y\geq a,\\
0\leq x\leq 2
\end{cases}
\]

表示的平面区域是一个三角形，则 $a$ 的取值范围是（　）

\begin{gkopts}(2)
\task （A）$a<5$
\task （B）$a\geq 7$
\task （C）$5\leq a<7$
\task （D）$a<5$ 或 $a\geq 7$
\end{gkopts}

\Needspace{4\baselineskip}
\textbf{解析}

在 $0\leq x\leq2$ 内，区域位于水平直线 $y=a$ 上方、直线 $y=x+5$ 下方。若 $a<5$，左右边界 $x=0$ 与 $x=2$ 上都有正长度竖边，形成四边形；若 $a=5$，左端收缩为一点，形成非退化三角形；若 $5<a<7$，两边界在 $x=a-5\in(0,2)$ 相交，也形成非退化三角形；$a=7$ 时只剩一个点，$a>7$ 时为空集。因此要求 $5\leq a<7$。

\gkanswer{答案：C，$5\leq a<7$}

\end{answer}
\begin{answer}
\gksolhead{7}{}

\textbf{题目}：

平面 $\alpha\parallel$ 平面 $\beta$ 的一个充分条件是（　）

（A）存在一条直线 $a$，$a\parallel\alpha$，$a\parallel\beta$

（B）存在一条直线 $a$，$a\subset\alpha$，$a\parallel\beta$

（C）存在两条平行直线 $a,b$，$a\subset\alpha$，$b\subset\beta$，$a\parallel\beta$，$b\parallel\alpha$

（D）存在两条异面直线 $a,b$，$a\subset\alpha$，$b\subset\beta$，$a\parallel\beta$，$b\parallel\alpha$

\Needspace{4\baselineskip}
\textbf{解析}

设 $\alpha$ 与 $\beta$ 相交，交线为 $l$。若 $a\subset\alpha$ 且 $a\parallel\beta$，则 $a$ 不与 $\beta$ 相交，所以 $a$ 在 $\alpha$ 内平行于 $l$；同理 $b\parallel l$。这与 $a,b$ 异面矛盾，故两平面不相交，从而平行。

\gkanswer{答案：D}

\end{answer}
\begin{answer}
\gksolhead{8}{}

\textbf{题目}：

对于函数① $f(x)=|x+2|$，② $f(x)=(x-2)^2$，③ $f(x)=\cos(x-2)$，判断如下两个命题的真假：

命题甲：$f(x+2)$ 是偶函数；

命题乙：$f(x)$ 在 $(-\infty,2)$ 上是减函数，在 $(2,+\infty)$ 上是增函数；

能使命题甲、乙均为真的所有函数的序号是（　）

\begin{gkopts}(4)
\task （A）①②
\task （B）①③
\task （C）②
\task （D）③
\end{gkopts}

\Needspace{4\baselineskip}
\textbf{解析}

① 有 $f(x+2)=|x+4|$，不是偶函数；③ 的平移为 $\cos x$，是偶函数，但原函数在题设的两个区间上不具有所要求的单调性。② 平移后为 $x^2$，而 $f(x)=(x-2)^2$ 在 $x<2$ 时递减、在 $x>2$ 时递增，故只有②同时满足条件。

\gkanswer{答案：C（②）}

\end{answer}
\gkpart{二、填空题}

\begin{answer}
\gksolhead{9}{}

\textbf{题目}：

$f'(x)$ 是 $f(x)=\dfrac{1}{3}x^3+2x+1$ 的导函数，则 $f'(-1)$ 的值是\gkblank{}。

\Needspace{4\baselineskip}
\textbf{解析}

求导得 $f'(x)=x^2+2$，代入 $x=-1$ 得 $f'(-1)=3$。

\gkanswer{答案：$3$}

\end{answer}
\begin{answer}
\gksolhead{10}{}

\textbf{题目}：

若数列 $\{a_n\}$ 的前 $n$ 项和 $S_n=n^2-10n\ (n=1,2,3,\ldots)$，则此数列的通项公式为\gkblank{}。

\Needspace{4\baselineskip}
\textbf{解析}

$a_1=S_1=-9$。当 $n\geq2$ 时，$a_n=S_n-S_{n-1}=n^2-10n-[(n-1)^2-10(n-1)]=2n-11$；该表达式在 $n=1$ 时也等于 $-9$，所以对所有正整数 $n$ 都成立。

\gkanswer{答案：$a_n=2n-11$}

\end{answer}
\begin{answer}
\gksolhead{11}{}

\textbf{题目}：

已知向量 $\boldsymbol a=(2,4)$，$\boldsymbol b=(1,1)$。若向量 $\boldsymbol b\perp(\boldsymbol a+\lambda\boldsymbol b)$，则实数 $\lambda$ 的值是\gkblank{}。

\Needspace{4\baselineskip}
\textbf{解析}

由垂直条件取点积：$(1,1)\cdot[(2,4)+\lambda(1,1)]=0$，即 $6+2\lambda=0$。

\gkanswer{答案：$\lambda=-3$}

\end{answer}
\begin{answer}
\gksolhead{12}{}

\textbf{题目}：

在 $\triangle ABC$ 中，若 $\tan A=\dfrac{1}{3}$，$C=150^\circ$，$BC=1$，则 $AB=$\gkblank{}。

\Needspace{4\baselineskip}
\textbf{解析}

$A$ 为三角形内角且 $\tan A=1/3$，故 $\sin A=1/\sqrt{10}$。由正弦定理，$AB=BC\,\sin C/\sin A=1\cdot(1/2)/(1/\sqrt{10})=\sqrt{10}/2$。

\gkanswer{答案：$\dfrac{\sqrt{10}}2$}

\end{answer}
\begin{answer}
\gksolhead{13}{}

\textbf{题目}：

2002 年在北京召开的国际数学家大会，会标是以我国古代数学家赵爽的弦图为基础设计的。弦图是由四个全等直角三角形与一个小正方形拼成的一个大正方形（如图）。如果小正方形的面积为 1，大正方形的面积为 25，直角三角形中较小的锐角为 $\theta$，那么 $\cos 2\theta$ 的值等于\gkblank{}。

〔图位：原卷赵爽弦图，待矢量重绘〕

\Needspace{4\baselineskip}
\textbf{解析}

设直角三角形较长、较短直角边为 $p,q$，斜边为 $c$。大正方形边长为斜边，故 $p^2+q^2=c^2=25$；中间小正方形边长为 $p-q$，所以 $(p-q)^2=1$。于是 $(p+q)^2=2(p^2+q^2)-(p-q)^2=49$，从而 $p+q=7$。较小角 $\theta$ 对着短边 $q$，故 $\cos2\theta=(p^2-q^2)/(p^2+q^2)=(p-q)(p+q)/25=7/25$。

\gkanswer{答案：$\dfrac7{25}$}

\end{answer}
\begin{answer}
\gksolhead{14}{}

\textbf{题目}：

已知函数 $f(x)$，$g(x)$ 分别由下表给出：

\begin{center}
\begin{tabular}{lrrr}
\toprule
$x$ & 1 & 2 & 3 \\
\midrule
$f(x)$ & 2 & 1 & 1 \\
\bottomrule
\end{tabular}
\end{center}

\begin{center}
\begin{tabular}{lrrr}
\toprule
$x$ & 1 & 2 & 3 \\
\midrule
$g(x)$ & 3 & 2 & 1 \\
\bottomrule
\end{tabular}
\end{center}

则 $f[g(1)]$ 的值为\gkblank{}；当 $g[f(x)]=2$ 时，$x=$\gkblank{}。

\Needspace{4\baselineskip}
\textbf{解析}

由表得 $g(1)=3$、$f(3)=1$，故 $f[g(1)]=1$。又 $g(y)=2$ 当且仅当 $y=2$，所以 $g[f(x)]=2$ 等价于 $f(x)=2$；查表只有 $x=1$。

\gkanswer{答案：$1$；$x=1$}

\end{answer}
\gkpart{三、解答题}

\begin{answer}
\gksolhead{15}{}

\textbf{题目}：

记关于 $x$ 的不等式 $\dfrac{x-a}{x+1}<0$ 的解集为 $P$，不等式 $|x-1|\leq 1$ 的解集为 $Q$。

(1) 若 $a=3$，求 $P$；

(2) 若 $Q\subseteq P$，求正数 $a$ 的取值范围。

\Needspace{4\baselineskip}
\textbf{解析}

\textbf{（1）} $a=3$ 时，临界点为 $x=-1,3$。分式在两根之间为负，且两端均不取，故解集为 $(-1,3)$。

\textbf{（2）} 对正数 $a$，分式的负值区间为 $P=(-1,a)$；而 $|x-1|\leq1$ 给出 $Q=[0,2]$。包含 $Q\subseteq P$ 等价于右端点 $2$ 严格小于 $a$，即 $a>2$。

\gkanswer{答案（1）：$P=(-1,3)$}

\gkanswer{答案（2）：$a>2$}

\end{answer}
\begin{answer}
\gksolhead{16}{}

\textbf{题目}：

数列 $\{a_n\}$ 中，$a_1=2$，$a_{n+1}=a_n+cn$（$c$ 是常数，$n=1,2,3,\ldots$），且 $a_1,a_2,a_3$ 成公比不为 1 的等比数列。

(1) 求 $c$ 的值；

(2) 求 $\{a_n\}$ 的通项公式。

\Needspace{4\baselineskip}
\textbf{解析}

\textbf{（1）} 由递推得 $a_2=2+c$、$a_3=2+3c$。三项成等比数列满足 $(2+c)^2=2(2+3c)$，化简为 $c(c-2)=0$。$c=0$ 时公比为 $1$，不合题意，故 $c=2$。

\textbf{（2）} 将递推从 $1$ 累加到 $n-1$：$a_n=a_1+2\sum_{j=1}^{n-1}j=2+n(n-1)=n^2-n+2$。

\gkanswer{答案（1）：$c=2$}

\gkanswer{答案（2）：$a_n=n^2-n+2$}

\end{answer}
\begin{answer}
\gksolhead{17}{}

\textbf{题目}：

如图，在 $\mathrm{Rt}\triangle AOB$ 中，$\angle OAB=\dfrac{\pi}{6}$，斜边 $AB=4$。$\mathrm{Rt}\triangle AOC$ 可以通过 $\mathrm{Rt}\triangle AOB$ 以直线 $AO$ 为轴旋转得到，且二面角 $B-AO-C$ 是直二面角。$D$ 是 $AB$ 的中点。

(1) 求证：平面 $COD\perp$ 平面 $AOB$；

(2) 求异面直线 $AO$ 与 $CD$ 所成角的大小。

〔图位：原卷立体几何图，待矢量重绘〕

\Needspace{4\baselineskip}
\textbf{解析}

\textbf{（1）} 二面角为直二面角，故平面 $AOC\perp$ 平面 $AOB$，交线为 $AO$。三角形 $AOC$ 在 $O$ 处为直角，故 $OC\perp AO$。在两个互相垂直的平面中，一平面内垂直交线的直线垂直于另一平面，因此 $OC\perp$ 平面 $AOB$。由于 $OC$ 位于平面 $COD$ 内，平面 $COD\perp$ 平面 $AOB$。

\textbf{（2）} 因 $AB=4$、$\angle OAB=\pi/6$ 且 $\triangle AOB$ 在 $O$ 处为直角，可取 $O=(0,0,0)$、$A=(2\sqrt3,0,0)$、$B=(0,2,0)$。旋转及直二面角条件允许取 $C=(0,0,2)$。中点 $D=(\sqrt3,1,0)$，于是 $\overrightarrow{CD}=(\sqrt3,1,-2)$。直线 $CD$ 与 $AO$ 的夹角取锐角，其余弦为 $\dfrac{|(\sqrt3,1,-2)\cdot(1,0,0)|}{\sqrt8}=\dfrac{\sqrt6}{4}$。

\gkanswer{答案（1）：见解析版}

\gkanswer{答案（2）：$\arccos\dfrac{\sqrt6}{4}$}

\end{answer}
\begin{answer}
\gksolhead{18}{}

\textbf{题目}：

某条公共汽车线路沿线共有 11 个车站（包括起点站和终点站）。在起点站开出的一辆公共汽车上有 6 位乘客，假设每位乘客在起点站之外的各个车站下车是等可能的。求：

(1) 这 6 位乘客在互不相同的车站下车的概率；

(2) 这 6 位乘客中恰有 3 人在终点站下车的概率。

\Needspace{4\baselineskip}
\textbf{解析}

\textbf{（1）} 11 个车站包括起点站，故每名乘客有 10 个可下车站。题面只规定每位乘客的边缘分布为均匀分布，未说明六位乘客的选择联合独立，所以本问的无条件概率不能由现有条件唯一确定。若附加相互独立假设，六人依次选择互不相同站点的概率为 $10\cdot9\cdot8\cdot7\cdot6\cdot5/10^6=189/1250$。

\textbf{（2）} 同理，题面没有给出六次选择间的联合关系，故无条件概率未定。在六人选择相互独立的额外假设下，指定三人下车于终点、另三人下车于其他站的概率为 $(1/10)^3(9/10)^3$；指定三人的选择有 $\binom63$ 种，因此所求概率为 $\binom63(1/10)^3(9/10)^3=729/50000$。

\gkanswer{答案（1）：U/PENDING；若进一步假设六位乘客的选择相互独立，则概率为 $\dfrac{10\cdot9\cdot8\cdot7\cdot6\cdot5}{10^6}=\dfrac{189}{1250}$。}

\gkanswer{答案（2）：U/PENDING；若进一步假设六位乘客的选择相互独立，则概率为 $\binom63(1/10)^3(9/10)^3=\dfrac{729}{50000}$。}

\end{answer}
\begin{answer}
\gksolhead{19}{}

\textbf{题目}：

如图，矩形 $ABCD$ 的两条对角线相交于点 $M(2,0)$，$AB$ 边所在直线的方程为 $x-3y-6=0$，点 $T(-1,1)$ 在 $AD$ 边所在直线上。

(1) 求 $AD$ 边所在直线的方程；

(2) 求矩形 $ABCD$ 外接圆的方程；

(3) 若动圆 $P$ 过点 $N(-2,0)$，且与矩形 $ABCD$ 的外接圆外切，求动圆 $P$ 的圆心的轨迹方程。

〔图位：原卷矩形坐标图，待矢量重绘〕

\Needspace{4\baselineskip}
\textbf{解析}

\textbf{（1）} 直线 $AB$ 的斜率为 $1/3$，因此 $AD$ 的斜率为 $-3$。过 $T(-1,1)$ 得 $y-1=-3(x+1)$，即 $AD:3x+y+2=0$。

\textbf{（2）} 联立 $AB:x-3y-6=0$ 与 $AD:3x+y+2=0$，得 $A=(0,-2)$。矩形对角线交点 $M=(2,0)$ 是外接圆圆心，半径平方为 $MA^2=(0-2)^2+(-2)^2=8$，故圆方程为 $(x-2)^2+y^2=8$。

\textbf{（3）} 设动圆圆心为 $P$，其半径为 $PN$；固定圆的圆心为 $M$、半径为 $2\sqrt2$。外切条件给出 $PM=PN+2\sqrt2$，即 $PM-PN=2\sqrt2$。两焦点为 $M=(2,0)$ 与 $N=(-2,0)$，故 $2a=2\sqrt2$、$c=2$、$b^2=c^2-a^2=2$。轨迹所在双曲线为 $x^2/2-y^2/2=1$。又 $PM>PN$，圆心更靠近 $N$，故取左支 $x\leq-\sqrt2$。

\gkanswer{答案（1）：$3x+y+2=0$}

\gkanswer{答案（2）：$(x-2)^2+y^2=8$}

\gkanswer{答案（3）：$\dfrac{x^2}{2}-\dfrac{y^2}{2}=1$ 的左支（$x\leq-\sqrt2$）}

\end{answer}
\begin{answer}
\gksolhead{20}{}

\textbf{题目}：

已知函数 $y=kx$ 与 $y=x^2+2\ (x\geq 0)$ 的图象相交于 $A(x_1,y_1)$，$B(x_2,y_2)$。$l_1,l_2$ 分别是 $y=x^2+2\ (x\geq 0)$ 的图象在 $A,B$ 两点的切线，$M,N$ 分别是 $l_1,l_2$ 与 $x$ 轴的交点。

(1) 求 $k$ 的取值范围；

(2) 设 $t$ 为点 $M$ 的横坐标，当 $x_1<x_2$ 时，写出 $t$ 以 $x_1$ 为自变量的函数式，并求其定义域和值域；

(3) 试比较 $|OM|$ 与 $|ON|$ 的大小，并说明理由（$O$ 是坐标原点）。

\begin{gktip}
\textbf{来源说明}：题面依据固定候选 PDF \texttt{普通高考/2007/2007北京文.pdf} 转录，source ID 为 \texttt{src-71817304e261efb7deaf51c2a5f400421a95ac91a90313d8f558bee701a624f9}，SHA-256 为 \texttt{e5356797259eb22cc2b4f23019ad7ab044f282bdfc2a65df94dbbdcc502cfff5}。PDF 只有题面，没有答案、解析或评分页；来源发布者、出版者、版次、复制链与官方身份未核，\texttt{official_answer} 保持 null。
\end{gktip}

\Needspace{4\baselineskip}
\textbf{解析}

\textbf{（1）} 交点横坐标满足 $x^2-kx+2=0$。两个不同交点要求判别式 $k^2-8>0$，且两根之和 $k>0$、积为 $2>0$，所以两个根均为正，得到 $k>2\sqrt2$。

\textbf{（2）} 在 $x=x_1$ 处抛物线切线斜率为 $2x_1$，切线方程为 $y=2x_1x-x_1^2+2$，故横轴交点横坐标 $t=x_1/2-1/x_1$。由根积 $x_1x_2=2$ 及 $0<x_1<x_2$ 得 $0<x_1<\sqrt2$。函数导数 $1/2+1/x_1^2>0$，且端点极限分别为 $-\infty$ 与 $0$，所以定义域为 $(0,\sqrt2)$，值域为 $(-\infty,0)$。

\textbf{（3）} 对 $x_2$ 处的切线同理，$N$ 的横坐标为 $x_2/2-1/x_2$。利用 $x_2=2/x_1$ 得 $x_2/2-1/x_2=1/x_1-x_1/2=-t$，故两横轴交点到原点的距离相同。

\gkanswer{答案（1）：$k>2\sqrt2$}

\gkanswer{答案（2）：$t=\dfrac{x_1}{2}-\dfrac1{x_1}$；定义域 $(0,\sqrt2)$，值域 $(-\infty,0)$}

\gkanswer{答案（3）：$|OM|=|ON|$}

\end{answer}